\SixthArticle{45}

Let $k$ be a positive integer (or also $=0$), and let us denote the first $k+1$ terms of the series $F(\alpha,\beta,\gamma,x)$ by $X$. The next term will be
\[
=\frac{\alpha.\alpha+1.\alpha+2\ .\ .\ .\ .\ \alpha+k.\beta.\beta+1.\beta+2\ .\ .\ .\ .\ \beta+k}
{1\ .\ 2\ .\ 3\ .\ .\ .\ .\ k+1.\gamma.\gamma+1.\gamma+2\ .\ .\ .\ .\ \gamma+k}x^{k+1}
\]
which can also be written as
\[
\frac{\Pi(\gamma-1)}{\Pi(\alpha-1)\Pi(\beta-1)}\ .\
\frac{\Pi(\alpha+k)\Pi(\beta+k)}{\Pi(k+1)\Pi(\gamma+k)}x^{k+1}
\]
and similarly for the following terms. Hence it follows that
\[
\mathrm{I.}\qquad
\frac{\Pi(\alpha+\beta-\gamma)(\Pi(-\gamma)}{\Pi(\alpha-\gamma)\Pi(\beta-\gamma)}\ .\ F(\alpha,\beta,\gamma,x)
\quad\text{can be expressed by}
\]
\[
\begin{aligned}&\frac{\Pi(\alpha+\beta-\gamma)\Pi(-\gamma)}{\Pi(\alpha-\gamma)\Pi(\beta-\gamma)}\ .\ X\\&+\frac{\Pi(\alpha+\beta-\gamma)\Pi(-\gamma)\Pi(\gamma-1)}{\Pi(\alpha-1)\Pi(\beta-1)\Pi(\alpha-\gamma)\Pi(\beta-\gamma)}\\&\qquad\Sigma\left\{\frac{\Pi(\alpha+k+t)\Pi(\beta+k+t)}{\Pi(k+t+1)\Pi(\gamma+k+t)}x^{k+1+t}\right\}\end{aligned}
\]
if $t$ is understood to take all the values $0,1,2,3$ etc. to infinity.

Similarly, $F(\alpha+1-\gamma,\beta+1-\gamma,2-\gamma,x)$ can be expressed by
\[
\frac{\Pi(1-\gamma)}{\Pi(\alpha-\gamma)\Pi(\beta-\gamma)}
\Sigma\left(\frac{\Pi(\alpha-\gamma+t)\Pi(\beta-\gamma+t)}{\Pi t\,\Pi(1-\gamma+t)}x^t\right)
\]
with $t$ determined as before; and since $\Pi(1-\gamma)=(1-\gamma)\Pi(-\gamma)$ and $\Pi(\gamma-1)=-(1-\gamma)\Pi(\gamma-2)$, it is evident that
\SourceLeaf{224}
\[
\begin{aligned}[c]
\mathrm{II.}\qquad&\frac{\Pi(\alpha+\beta-\gamma)\Pi(\gamma-2)}{\Pi(\alpha-1)\Pi(\beta-1)}x^{1-\gamma}F(\alpha+1-\gamma,\beta+1-\gamma,2-\gamma,x)\\
&=-\frac{\Pi(\alpha+\beta-\gamma)\Pi(-\gamma)\Pi(\gamma-1)}{\Pi(\alpha-1)\Pi(\beta-1)\Pi(\alpha-\gamma)\Pi(\beta-\gamma)}
\\&\qquad\Sigma\left\{\frac{\Pi(\alpha-\gamma+t)\Pi(\beta-\gamma+t)}{\Pi t\,\Pi(1-\gamma+t)}x^{1+t-\gamma}\right\}
\end{aligned}
\]
Hence formula 86 can also be written as follows:
\[
\begin{aligned}[c]
&F(\alpha,\beta,\alpha+\beta+1-\gamma,1-x)\\
&=\frac{\Pi(\alpha+\beta-\gamma)\Pi(-\gamma)}{\Pi(\alpha-\gamma)\Pi(\beta-\gamma)}X\\
&\quad+\frac{\Pi(\alpha+\beta-\gamma)\Pi(-\gamma)\Pi(\gamma-1)}{\Pi(\alpha-1)\Pi(\beta-1)\Pi(\alpha-\gamma)\Pi(\beta-\gamma)}
\\&\qquad\Sigma\Bigl\{\frac{\Pi(\alpha+k+t)\Pi(\beta+k+t)}{\Pi(k+t+1)\Pi(\gamma+k+t)}x^{k+1+t}
\\&\qquad-\frac{\Pi(\alpha-\gamma+t)\Pi(\beta-\gamma+t)}{\Pi(t-\gamma+1)\Pi t}x^{1+t-\gamma}\Bigr\}
\end{aligned}
\]
This expression immediately shows that each difference under the sign $\Sigma$ becomes $=0$ if $\gamma=-k$ is assumed. But since $\Pi(\gamma-1)$ also becomes an infinitely large quantity here, it is evident that the product can be finite. To express its value by finite quantities, let us first set $\gamma+k=\omega$, which gives
\[
\Pi(\gamma-1).\gamma.(\gamma+1)(\gamma+2)\ .\ .\ .\ .\ (\gamma+k-1)\omega=\Pi\omega
\]
or
\[
\Pi(\gamma-1)=\frac{\Pi\omega}{\omega(\omega-1)(\omega-2)\ .\ .\ .\ .\ (\omega-k)}
\]
\Needspace{15\baselineskip}
The essential point is therefore for us to see what
\[
\begin{aligned}&\frac{1}{\omega}\Bigl\{\frac{\Pi(\alpha-\gamma+t+\omega)\Pi(\beta-\gamma+t+\omega)}{\Pi(t-\gamma+1+\omega)\Pi(t+\omega)}x^{1+t-\gamma+\omega}\\&\qquad-\frac{\Pi(\alpha-\gamma+t)\Pi(\beta-\gamma+t)}{\Pi(t-\gamma+1)\Pi t}x^{1+t-\gamma}\Bigr\}\end{aligned}
\]
becomes as $\omega$ decreases towards zero. By known principles, however, this gives
\[
-\frac{\mathrm{d}U}{\mathrm{d}\gamma}
\]
if, for brevity, we set
\[
\frac{\Pi(\alpha-\gamma+t)\Pi(\beta-\gamma+t)}{\Pi(t-\gamma+1)\Pi(t-k-\gamma)}x^{1+t-\gamma}=U
\]
and regard $\gamma$ alone as variable. But this gives
\[
\begin{aligned}&\frac{\mathrm{d}U}{U\mathrm{d}\gamma}
=-\Psi(\alpha-\gamma+t)-\Psi(\beta-\gamma+t)\\&\qquad+\Psi(t-\gamma+1)+\Psi(t-k-\gamma)-\log x\end{aligned}
\]
Hence, for $\gamma=-k$, it follows that
\SourceLeaf{225}
\[
\begin{aligned}[c]
[88]\quad &F(\alpha,\beta,\alpha+\beta+1+k,1-x)\\
&=\frac{\Pi(\alpha+\beta+k)\Pi k}{\Pi(\alpha+k)\Pi(\beta+k)}X\\
&\quad+\frac{\Pi(\alpha+\beta+k)\Pi k}{\Pi(\alpha-1)\Pi(\beta-1)\Pi(\alpha+k)\Pi(\beta+k)(-1)(-2)\ldotp\ldotp\ldotp\ldotp(-k)}\\&\qquad\Sigma\{(\log x\\
&\qquad+\Psi(\alpha+t+k)+\Psi(\beta+t+k)-\Psi(t+k+1)-\Psi t)
\\&\qquad\frac{\Pi(\alpha+t+k)\Pi(\beta+t+k)}{\Pi(t+k+1)\Pi t}x^{1+t+k}\}\\
&=\frac{\Pi(\alpha+\beta+k)\Pi k}{\Pi(\alpha+k)\Pi(\beta+k)}X
\pm\frac{\Pi(\alpha+\beta+k)x^{1+k}}{\Pi(\alpha-1)\Pi(\beta-1)\Pi(k+1)}Y
\end{aligned}
\]
where
\[
\begin{aligned}[c]
&Y=\{\log x+\Psi(\alpha+k)+\Psi(\beta+k)-\Psi(k+1)-\Psi(0)\}\\&\qquad F(\alpha+k+1,\beta+k+1,k+2,x)\\
&+A\frac{\alpha+k+1.\beta+k+1}{1. k+2}x\\
&+(A+B)\frac{\alpha+k+1.\alpha+k+2.\beta+k+1.\beta+k+2}{1.2. k+2. k+3}xx\\
&+(A+B+C)\\&\qquad\frac{\alpha+k+1.\alpha+k+2.\alpha+k+3.\beta+k+1.\beta+k+2.\beta+k+3}{1.2.3. k+2. k+3. k+4}x^3\\
&+\text{etc.}
\end{aligned}
\]
and
\[
\begin{aligned}[c]
A&=\frac{1}{\alpha+k+1}+\frac{1}{\beta+k+1}-\frac{1}{k+2}-1\\
B&=\frac{1}{\alpha+k+2}+\frac{1}{\beta+k+2}-\frac{1}{k+3}-\frac12\\
C&=\frac{1}{\alpha+k+3}+\frac{1}{\beta+k+3}-\frac{1}{k+4}-\frac13,\\
&\text{etc.}
\end{aligned}
\]
and the upper or lower sign is to be taken according as $k$ is an even or an odd number.

\SixthArticle{46}
Thus, this is how $F(\alpha,\beta,\alpha+\beta+1-\gamma,1-x)$ is transformed if $\gamma$ is $0$ or a negative integer. We can treat the case $\gamma=+1$ in exactly the same way; or, more briefly, we can set $k=-1$ in the preceding reasoning, whereby $X$ disappears entirely and we obtain:
\SourceLeaf{226}
\[
\begin{aligned}[c]
[89]\quad &F(\alpha,\beta,\alpha+\beta,1-x)\\
&=-\frac{\Pi(\alpha+\beta-1)}{\Pi(\alpha-1)\Pi(\beta-1)}
\\&\qquad\{\log x+\Psi(\alpha-1)+\Psi(\beta-1)-2\Psi(0)\}F(\alpha,\beta,1,x)\\
&\quad-\frac{\Pi(\alpha+\beta-1)}{\Pi(\alpha-1)\Pi(\beta-1)}
\{A\frac{\alpha.\beta}{1}x\\
&\quad+(A+B)\frac{\alpha.\alpha+1.\beta.\beta+1}{1.2.1.2}xx\\
&\quad+(A+B+C)\frac{\alpha.\alpha+1.\alpha+2.\beta.\beta+1.\beta+2}{1.2.3.1.2.3}x^3\\
&\quad+\text{etc.}\}
\end{aligned}
\]
where
\[
\begin{aligned}&A=\frac1\alpha+\frac1\beta-2,\\&
B=\frac1{\alpha+1}+\frac1{\beta+1}-\frac22.\\&
C=\frac1{\alpha+2}+\frac1{\beta+2}-\frac23,\\&\text{etc.}\end{aligned}
\]
Thus, for example, for $\alpha=\frac12$, $\beta=\frac12$, we obtain (compare formulas 52, 71)
\[
\begin{aligned}[c]
[90]\quad &F(\tfrac12,\tfrac12,1,1-x)\\
&=-\frac1\pi\log\tfrac1{16}x. F(\tfrac12,\tfrac12,1,x)\\
&\quad-\frac1\pi\{2.\frac{1.1}{2.2}x
\\&\qquad+(2+\tfrac13)\frac{1.1.3.3}{2.2.4.4}xx
\\&\qquad+(2+\tfrac13+\tfrac2{15})\frac{1.1.3.3.5.5}{2.2.4.4.6.6}x^3\\
&\qquad+(2+\tfrac4{3.4}+\tfrac4{5.6}+\tfrac4{7.8})
\frac{1.1.3.3.5.5.7.7}{2.2.4.4.6.6.8.8}x^4+\text{etc.}\}\\
&=-\frac1\pi\{\log\tfrac1{16}xF(\tfrac12,\tfrac12,1,x)
\\&\qquad+\tfrac12x+\tfrac{21}{64}xx+\tfrac{185}{768}x^3+\tfrac{18655}{98304}x^4\\
&\qquad+\tfrac{102501}{655360}x^5+\tfrac{1394239}{10485760}x^6+\text{etc.}\}
\end{aligned}
\]
Finally, the third case, where $\gamma$ is a positive integer greater than unity, need not be treated separately, since
\[
\begin{aligned}&F(\alpha,\beta,\alpha+\beta+1-\gamma,1-x)\\&=x^{\gamma-1}F(\alpha+1-\gamma,\beta+1-\gamma,\alpha+\beta+1-\gamma,1-x)\end{aligned}
\]
and the transformation of the series $F(\alpha+1-\gamma,\beta+1-\gamma,\alpha+\beta+1-\gamma,1-x)$ for $\gamma>1$ reduces of itself to the first case.

\Needspace{18\baselineskip}
\SixthArticle{47}
We turn to other transformations, among which the substitution $x=\frac{y}{y-1}$ shall take first place. This gives $\mathrm{d}x=-\frac{\mathrm{d}y}{(y-1)^2}$, and therefore
\[
\frac{\mathrm{d}P}{\mathrm{d}x}=-\frac{\mathrm{d}P}{\mathrm{d}y}(1-y)^2,
\quad\text{differentiating again gives}
\]
\[
\mathrm{d}\frac{\mathrm{d}P}{\mathrm{d}x}=-(1-y)^2\mathrm{d}\frac{\mathrm{d}P}{\mathrm{d}y}
+2(1-y)\mathrm{d}P,\quad\text{therefore}
\]
\[
\frac{\mathrm{dd}P}{\mathrm{d}x^2}=+(1-y)^4\frac{\mathrm{dd}P}{\mathrm{d}y^2}
-2(1-y)^3\frac{\mathrm{d}P}{\mathrm{d}y}
\]
\SourceLeaf{227}
On substituting these values, equation 80 becomes
\[
\begin{aligned}&0=\alpha\beta P+(1-y)(\gamma+(\alpha+\beta-1-\gamma)y)\frac{\mathrm{d}P}{\mathrm{d}y}\\&+(1-y)(y-yy)\frac{\mathrm{dd}P}{\mathrm{d}y^2}\end{aligned}
\]
But to obtain an equation similar to 80, let us set $P=(1-y)^\mu P'$, which gives
\[
\begin{aligned}[c]
\frac{\mathrm{d}P}{\mathrm{d}y}
&=-\mu(1-y)^{\mu-1}P'+(1-y)^\mu\frac{\mathrm{d}P'}{\mathrm{d}y}\\
\frac{\mathrm{dd}P}{\mathrm{d}y^2}
&=(\mu\mu-\mu)(1-y)^{\mu-2}P'-2\mu(1-y)^{\mu-1}\frac{\mathrm{d}P'}{\mathrm{d}y}
\\&\qquad+(1-y)^\mu\frac{\mathrm{dd}P'}{\mathrm{d}y^2}
\end{aligned}
\]
On substituting these expressions and dividing by $(1-y)^\mu$, this gives
\[
\begin{aligned}[c]
&0=P'\{\alpha\beta-\mu(\gamma+(\alpha+\beta-1-\gamma)y+y(\mu\mu-\mu)\}\\
&+\frac{\mathrm{d}P'}{\mathrm{d}y}\{\gamma+(\alpha+\beta-1-\gamma)y-2\mu y\}(1-y)\\
&+\frac{\mathrm{dd}P'}{\mathrm{d}y^2}\{y-yy\}(1-y)
\end{aligned}
\]
Let us determine $\mu$ so that the multiplier of $P'$ becomes divisible by $1-y$. This is achieved by setting either $\mu=\alpha$ or $\mu=\beta$. The former value changes the preceding equation into
\[
\begin{aligned}&0=\alpha(\beta-\gamma)P'+(\gamma-(\gamma+\alpha+1-\beta)y)\frac{\mathrm{d}P'}{\mathrm{d}y}\\&+(y-yy)\frac{\mathrm{dd}P'}{\mathrm{d}y^2}\end{aligned}
\]
or
\[
\begin{aligned}&0=\alpha(\gamma-\beta)P'-(\gamma-(\gamma-\beta+\alpha+1)y)\frac{\mathrm{d}P'}{\mathrm{d}y}\\&-(y-yy)\frac{\mathrm{dd}P'}{\mathrm{d}y^2}\end{aligned}
\]
which must be satisfied so that, for $y=0$, $P'=1$ and $\frac{\mathrm{d}P'}{\mathrm{d}y}=\frac{\alpha(\gamma-\beta)}{\gamma}$. From this, however, $P'=F(\alpha,\gamma-\beta,\gamma,y)$ is deduced, and therefore
\[
\begin{aligned}&[91]\quad F(\alpha,\beta,\gamma,x)=(1-y)^\alpha F(\alpha,\gamma-\beta,\gamma,y)\\&=(1-x)^{-\alpha}F(\alpha,\gamma-\beta,\gamma,-\tfrac{x}{1-x})\end{aligned}
\]
Had we adopted for $\mu$ the other value $\beta$, exactly the same procedure would have given
\[
[92]\quad F(\alpha,\beta,\gamma,x)=(1-x)^{-\beta}F(\beta,\gamma-\alpha,\gamma,-\tfrac{x}{1-x})
\]
\Needspace{11\baselineskip}
which formula also follows of itself from the preceding one merely by interchanging the elements $\alpha$, $\beta$. With the aid of the formula just found, the values of our series for negative values of the fourth element are always reduced to the values of similar series for positive values of the fourth element lying between $0$ and $1$, since
\[
F(\alpha,\beta,\gamma,-x)=(1+x)^{-\alpha}F(\alpha,\gamma-\beta,\gamma,\tfrac{x}{1+x})
\]
\SourceLeaf{228}
